\chapter*{Abstract}
\phantomsection
\addcontentsline{toc}{chapter}{Abstract}

The massive adoption of cloud computing has transformed the management of technology infrastructure, but it has also introduced a persistent problem: a substantial share of cloud spending is wasted on underused or oversized resources, around one third according to industry studies, and manually detecting these inefficiencies is increasingly difficult in dynamic and complex architectures. This work describes the design, development, and validation of InfraEye, a platform for cloud infrastructure optimization in the Huawei Cloud ecosystem. The solution collects utilization metrics directly from the provider's APIs, without installing third-party agents, and combines two complementary mechanisms: a recommendation engine based on operational threshold rules that evaluates every instance, and the unsupervised Isolation Forest algorithm, which provides an additional multivariate classification of behaviors that deviate from the reference pattern. Based on these detections, the platform generates explainable recommendations and produces Terraform code that materializes the proposed changes, at all times subject to user review and approval.

Requirements were grounded in primary research, consisting of a survey with 162 responses and three interviews with specialists, which confirmed the problem and defined design criteria centered on recommendation transparency and human oversight of changes. Validation comprised a suite of 317 automated tests, an evaluation of Isolation Forest on a 400-machine test sample from a public cluster trace dataset against an expert criterion based on mean CPU utilization, and the functional verification of the complete analysis circuit in a controlled Huawei Cloud environment. On the test sample, the model achieved 0.980 accuracy, 0.795 precision, 1.000 recall, and an F1-score of 0.886; it detected all 31 positive cases defined by the expert criterion and produced 8 false positives. Measuring the actual economic savings is left as future work. Finally, the work assesses the commercial viability of the proposal through a SaaS model with a free tier and tiered plans, calibrated with prices and conversion rates from active market solutions: with a 17\% discount rate, all three projected scenarios yield a positive net present value, with a 30.1\% IRR in the neutral scenario and a minimal margin in the unfavorable trajectory, supporting the technical and economic viability of the platform, conditional on free-tier conversion.

\vspace{1cm}

\noindent\textbf{Keywords:} cloud computing, cost optimization, FinOps, anomaly detection, Isolation Forest, Infrastructure as Code, Terraform, Huawei Cloud, unsupervised machine learning.
